\section{Task 3: Kakuro as a Constraint Satisfaction Problem}

\subsection{Formulation}
We model Kakuro as a CSP. Each white cell $(r,c)$ is a variable with domain $\{1,\dots,9\}$. A \emph{run} is a maximal sequence of consecutive white cells in a row or column following a clue cell. Every run $R$ with clue $s_R$ imposes two $n$-ary constraints: (i)~all-different, requiring pairwise distinct values in $R$, and (ii)~the values in $R$ must sum to $s_R$. Every white cell belongs to exactly two runs (one horizontal and one vertical). The $5\times5$ example has 12 variables and 8 runs (clue sum 60), whereas the $9\times9$ test case has 48 variables and 28 runs, with horizontal and vertical clue sums both equal to 249.

\subsection{Approach and Visualization}
The solver combines constraint propagation with backtracking search. Valid digit combinations for each run length and target sum are precomputed and cached. During propagation, combinations inconsistent with already assigned values or the current domains of unassigned cells in a run are discarded: every unassigned cell must share at least one candidate with the remaining digits of a combination, and every remaining digit must appear in at least one unassigned cell's domain. The domain of each unassigned cell is then intersected with the union of digits in the surviving combinations minus the digits already assigned in that run. Because narrowing a cell's domain can eliminate combinations in its intersecting perpendicular run, propagation iterates to a fixed point, triggering backtracking whenever a domain or combination set becomes empty. Search supports Minimum Remaining Values (MRV), Degree, MRV with Degree tie-breaking, and Least Constraining Value (LCV), saving and restoring domains and assignments at each branching step.

To visualize the solving process without modifying the core solver, an event bus (\texttt{EventEmitter}) records solver events and snapshots for a Tkinter GUI. The interface renders candidate domains inside each unassigned cell as a $3\times3$ subgrid, highlights constraint propagation (\texttt{DOMAIN\_UPDATE}, yellow), value assignment (\texttt{ASSIGN}, green), and backtracking (\texttt{BACKTRACK}, red), and provides puzzle/heuristic selection alongside step-by-step and variable-speed playback.

\subsection{Experiments}
All benchmark configurations were evaluated on the $9\times9$ board with propagation enabled and a cap of 50{,}000 assignments (Table~\ref{tab:kakuro}).

\begin{table}[htbp]
\caption{Heuristic configurations on the $9\times9$ Kakuro board.}
\label{tab:kakuro}
\centering
\begin{tabular}{llrrr}
\hline
Variable & Value & Assign. & Backtr. & Time (s) \\
\hline
none       & none & 77       & 29       & 0.0456 \\
none       & LCV  & 50       & 2        & 0.0266 \\
MRV        & none & 64       & 16       & 0.0350 \\
MRV        & LCV  & 62       & 14       & 0.0332 \\
MRV+Degree & none & 49       & 1        & 0.0331 \\
MRV+Degree & LCV  & 52       & 4        & 0.0368 \\
Degree     & none & 8{,}269  & 8{,}221  & 3.7506 \\
Degree     & LCV  & 44{,}965 & 44{,}917 & 22.1582 \\
\hline
\end{tabular}
\end{table}

With fixed-point domain-aware propagation pruning combinations across intersecting runs, all eight configurations complete within the 50{,}000-assignment cap. \emph{MRV+Degree} without LCV achieves the fewest assignments (49 assignments, 1 backtrack, 0.0331~s), followed closely by \emph{none+LCV} (50 assignments, 2 backtracks, 0.0266~s) and \emph{MRV+Degree+LCV} (52 assignments, 4 backtracks, 0.0368~s). On this board, LCV reduces search under fixed ordering (77 to 50) and slightly improves standalone MRV (64 to 62), whereas it slightly increases assignments for MRV+Degree (49 to 52) and substantially degrades standalone Degree (8{,}269 to 44{,}965). Overall, Degree alone remains orders of magnitude slower than MRV-based selection (3.7506~s and 22.1582~s) because it ignores propagated domain sizes. Disabling propagation increases the search to 186{,}504 assignments (186{,}456 backtracks, over $2{,}900\times$ more than MRV with propagation). Across 20 repeated runs, the solver returned identical valid solutions and exact search counts (12 assignments and 0 backtracks on $5\times5$; 64 assignments and 16 backtracks on $9\times9$), properly rejected unsolvable boards, and solved a multi-solution test board. An independent exhaustive counter confirmed that both the $5\times5$ and $9\times9$ assignment boards have unique solutions.